% Bagian: second-order-logic
% Bab: sol-and-sets
% Subbagian: introduction

\documentclass[../../../include/open-logic-section]{subfiles}

\begin{document}

\olfileid[id]{sol}{set}{int}

\olsection{Pendahuluan}

Karena logika orde kedua dapat menguantifikasi atas himpunan-himpunan bagian
domain maupun atas fungsi, dapat diduga bahwa setidaknya sebagian teori
himpunan dapat dijalankan dalam logika orde kedua. Dengan ``dijalankan,''
yang kita maksud ialah bahwa sifat dan pernyataan teori himpunan dapat
diungkapkan dalam logika orde kedua, dan hal itu dapat dilakukan tanpa
kosakata nonlogis khusus untuk himpunan (misalnya, !!{predicate} keanggotaan
dalam teori himpunan). Sebagai contoh, kita dapat mendefinisikan gabungan dan
irisan himpunan serta relasi himpunan bagian, tetapi juga membandingkan ukuran
himpunan dan menyatakan hasil seperti Teorema Cantor.

\end{document}
